\section{El residuo de la raíz de dos fibras}
Consideremos una fibra de fila $i$ y una fibra de columna $h$ en un término
del exterior común. Sea $Y_1,Y_2$ su ocupación, cada una con valores en $\{0,1\}$.
A campo basal congelado escribimos $m_i=\mathbb E Y_i$ ($i=1,2$) y
$\Sigma_{ih}=\operatorname{Cov}(Y_i,Y_h)$ para los índices locales $i,h\in\{1,2\}$.
Definimos $f(y)=y^2-2y$, $l(y)=-y$ y
\[
 P=f_1f_2,\qquad B=P-l_1f_2-f_1l_2.
\]
Aquí los subíndices indican el argumento y $B$ designa el polinomio del prefactor.
Para ocupaciones binarias $f(Y_i)=l(Y_i)$, luego $(P+B)(Y)=0$ en cada átomo.

\paragraph{La identidad con covarianza.}
Sea $v_i=m_i(1-m_i)$. Así $l(m_i)=f(m_i)+v_i$ y
$(P+B)(m)=-v_1 f(m_2)-v_2 f(m_1)$. Puesto que $f''=2$, se obtiene
\[
 (P+B)(m)+\tfrac12\operatorname{tr}(\Sigma P''(m))
   = f'(m_1)f'(m_2)\Sigma_{12}
   =4(1-m_1)(1-m_2)\Sigma_{12}.
\]
Las contribuciones diagonales de la traza cancelan exactamente los términos
con $v_1,v_2$. Sólo la covarianza entre las dos fibras queda en la raíz.

\paragraph{Por qué comparten su amplitud en la máscara.}
En las notaciones de \emph{Media matriz}, partes 04 y 05, sea
$E_{ab}=\int z_x^2\,dM_a^R\,dM_b^C$. La acción aporta $E_{22}/2$
y el prefactor aporta $(E_{22}-E_{12}-E_{21})/2$. El núcleo $z_x$ se fija
por las mismas medidas $M_1$ en las tres integrales. Extraer la pareja de átomos
de grupos $g,g'$ da el mismo factor $h_g h_{g'}z_x(r_i,c_h)^2/2$ delante
de $P$ y de $B$, respectivamente. Se extrae el mismo coeficiente de Taylor
del núcleo cuando se desarrollan las perturbaciones de $M_1$.

\paragraph{Aplicación al soporte máximo.}
En coste $l$, las dos marcas dejan a lo sumo $l-2$ factores lineales adicionales.
Para ocupar $l$ grupos distintos, las dos fibras y todos estos factores deben
pertenecer a grupos distintos. En el bosque basal la ley se factoriza por grupos,
incluso al variar sus campos locales. Por ello $\Sigma_{12}=0$, los factores
privados sólo aportan sus medias y la misma identidad vale en esos campos.
En la traza, el segundo diferencial de cada factor privado es cero y los
bloques de covarianza entre grupos distintos son cero. La reducción de ramas
de \emph{Patrones finitos}, parte 04, permite derivar la identidad así obtenida.
Si las dos fibras pertenecen a un solo grupo, el número $h_b$ de grupos de
este término es a lo sumo $l-1$. Su déficit $l+1-h_b$ es al menos dos. La cuenta
de incidencias $d=w+1-\sum_b(l_b+1-h_b)-L_I$, con $L_I\ge0$, da $d\le w-1$.
La conclusión concierne a esta raíz y a sus decoraciones de soporte máximo.
